% Sixfold contribution, 21 November 2026.
% Author: GPT-6 Astra Pro, /root/sixfolds; inherited effort not exposed.
% Parent supplies the preamble and shared bibliography.
\section{Cyclic sixfold sources, exact period rank, and small detector bounds}
\label{sec:HD-S}

\paragraph{Scope.}
The cyclic genus-four source reaches the split Gaussian and Eisenstein
rational Hermitian spaces. The primitive period maps for character orders
\(3,4,6\) all have exact differential rank nine. The principal Gaussian
source has the even integral lattice \(U^3\). A cubic Eisenstein source
has an explicit principal quotient of lattice \(U^3\), integrally isometric
to the balanced product lattice \(I_{3,3}\).
The direct symmetric-product construction below gives
\[
 D^{\rm src}_{6,i}=746496,\qquad D^{\rm src}_{6,\omega}=30375.
\]
The ordinary determinant convention here is
\(\det(h)\in\mathbb Q^\times/N_{K/\mathbb Q}K^\times\);
the split class in rank six is \([-1]\).
The full list of omissions for this specified construction is proved below.

\subsection{Integral lattices and the source image}

\paragraph{HD-S01 (an integral topological decomposition).}
For a connected etale cyclic \(d\)-cover \(C\to X\), with \(g(X)=g\),
there is an orthogonal symplectic decomposition
\[
 H_1(C,\mathbb Z)=M_0\ \perp\
 \bigoplus_{j=2}^{g}
  \bigl(\mathbb Z[C_d]a_j\oplus\mathbb Z[C_d]b_j\bigr),
\]
where \(M_0\) has rank two and trivial deck action, and
\[
 E(a_{j,r},b_{k,s})=\delta_{jk}\delta_{rs}.
\]
To prove it, choose a symplectic marking placing the surjective cyclic
monodromy on the first handle. Cut \(X\) along a separating circle into
a once-punctured torus and a genus-\((g-1)\) surface with one boundary.
The latter lifts to \(d\) copies; their handle cycles give the displayed
regular modules and pairings. The former lifts to a genus-one surface
with \(d\) boundaries. After the lifted complementary surfaces are glued,
its boundary classes vanish, leaving the rank-two homology of its capped
torus. The capped covering is an etale torus cover, and its deck
translations are trivial on homology. Mayer--Vietoris gives the integral
direct sum, and disjoint supporting subsurfaces give orthogonality.
The initial marking follows by taking a primitive integral lift of the
cyclic cohomology character and extending its dual primitive vector to a
symplectic basis.

\paragraph{HD-S02 (primitive polarizations).}
For \(d=4\), the primitive factor
\(B=\ker(1+\sigma^2)^0\) has on each regular handle the saturated lattice
\[
 A=a_0-a_2,\quad iA=a_1-a_3,\quad
 B_0=b_0-b_2,\quad iB_0=b_1-b_3.
\]
Its only nonzero independent alternating pairings are
\(E(A,B_0)=E(iA,iB_0)=2\).
Consequently
\[
 \Theta_C|_B=2\Xi,\qquad \Xi\text{ principal},\qquad
 \ker(\Theta_C|_B)=B[2].
\]
For \(g=4\), the restricted type is \((2^6)\), and its kernel has order
\(2^{12}\); for \(g=3\), the type is \((2^4)\), with order \(2^8\).
The integral Gaussian Hermitian block of \(\Xi\) is
\(\begin{psmallmatrix}0&i\\-i&0\end{psmallmatrix}\), isometric by a
diagonal unit change to \(U=\begin{psmallmatrix}0&1\\1&0\end{psmallmatrix}\).
Thus the genus-four principal source has lattice \(U^3\).
The explicit quarter-turn deck action and these pairings prove the assertion;
no inference is made just from the order of the polarization kernel.

For \(d=3\), the primitive factor is the full norm Prym \(P\).
Its augmentation pairing on each handle has Smith invariants \(1,3\),
giving type \((1^3,3^3)\).
Write \(\omega=\sigma\), \(\delta_3=\sqrt{-3}=1+2\omega\), and choose
\(A=a_0-a_1,\ B_0=b_0-b_1\).
For the convention
\[
 h(x,y)=E(x,\omega y)-\bar\omega E(x,y)
\]
(conjugate-linear in the first argument), the Hermitian block is
\[
 \begin{pmatrix}0&\delta_3\\-\delta_3&0\end{pmatrix}.
\]
Replacing the lattice by the index-three overlattice with basis
\(A,B_0/\delta_3\) gives Gram matrix \(U\).
In three blocks this constructs a degree-27 quotient
\(P\to P'\) with a principal polarization \(\Xi'\) satisfying
\(\phi^*\Xi'=\Theta_C|_P\), and principal lattice \(U^3\).

For \(d=6\), put
\[
 S(t)=\frac{t^6-1}{t^2-t+1}=t^4+t^3-t-1.
\]
The primitive integral part of each regular module is generated by
\(S(\sigma)a,\sigma S(\sigma)a\), and similarly for \(b\).
This follows by monic polynomial division in
\(\mathbb Z[t]/(t^6-1)\). The cross-pairing matrix is
\(\begin{psmallmatrix}4&2\\2&4\end{psmallmatrix}\), of Smith invariants \(2,6\).
Thus \(\Theta_C|_B\) has type \((2^3,6^3)\).
Its half has type \((1^3,3^3)\); with \(\omega=\sigma-1\), the Hermitian
block before taking the half is
\(\begin{psmallmatrix}0&2\delta_3\\-2\delta_3&0\end{psmallmatrix}\).
The preceding index-27 quotient therefore gives another explicit
principal \(U^3\) target. This compares lattices by actual quotient maps,
without identifying the order-three and order-six cover moduli.

\paragraph{HD-S03 (fields and ordinary determinant).}
The primitive cyclotomic factor of character order \(e>2\) for a
genus-four etale cyclic cover has abelian dimension \(3\varphi(e)\).
Indeed, the regular-handle decomposition gives six copies of each
nontrivial complex character in \(H^1\), three holomorphic and three
antiholomorphic. A sixfold factor of this specified primitive construction
therefore requires \(\varphi(e)=2\), namely \(e=3,4,6\).
Its field is \(\mathbb Q(\omega)\) or \(\mathbb Q(i)\).

Projecting the \(a\)-cycles and \(b\)-cycles of HD-S01 onto the primitive
factor gives two rational totally isotropic K-subspaces, each of dimension
three, paired perfectly. The Hermitian space is therefore split:
in dual bases its Gram matrix is \(U^3\), of determinant \(-1\).
A K-linear rational change of basis \(F\) changes its determinant by
\(N(\det F)\). A polarization similitude of rational multiplier \(c\)
changes it by \(c^6=N(c^3)\). Thus every polarized K-linear transport
of this source has ordinary determinant class \([-1]\).
Koike's published p.567 writes his determinant quotient using
\((K^\times)^2\) and labels his family by determinant one; the
explicit split-lattice calculation above specifies the norm-quotient
convention used here. The signed determinant is \(+1\).

\paragraph{HD-S03a (exhaustive split determinant criterion).}
For every imaginary quadratic field \(K\), a nondegenerate Hermitian
space of signature \((3,3)\) has ordinary determinant class \([-1]\)
exactly when it is \(K\)-isometric to \(U^3\).
Here is a direct proof of the converse, including the actual changes
of basis. The rational quadratic form \(q(v)=h(v,v)\) on the underlying
twelve-dimensional \(\mathbb Q\)-space is indefinite. Meyer's theorem
gives a nonzero \(v\) with \(h(v,v)=0\).
The precise published input is J.-P. Serre,
\emph{A Course in Arithmetic}, Chapter IV, \S3.2, Corollary 2 to
Theorem 8, printed p.43: an indefinite rational quadratic form of
rank at least five has a nontrivial zero.
Choose \(w\) with \(h(v,w)=1\). With the convention conjugate-linear
in the first argument, replacing \(w\) by
\(w-h(w,w)v/2\) makes it isotropic and preserves \(h(v,w)=1\).
These two vectors have Gram matrix \(U\), and their orthogonal
complement has signature \((2,2)\). Its underlying rational quadratic
form has rank eight, so the same construction splits a second \(U\).
The remaining Hermitian plane has signature \((1,1)\). Diagonalizing
it gives \(\langle a,b\rangle\), with \(a,b\in\mathbb Q^\times\).
The determinant hypothesis and the two removed \(U\)-blocks give
\(ab=-N_{K/\mathbb Q}(c)\) for some \(c\in K^\times\).
The vector \((c/a,1)\) has norm
\(aN(c/a)+b=N(c)/a+b=0\). Applying the displayed \(v,w\)
construction once more makes this plane \(U\).
Conversely, \(\det(U^3)=-1\), and a \(K\)-linear basis change
multiplies a determinant by a field norm, as already proved.
Thus every signature-\((3,3)\) target in the determinant class
\([-1]\) is covered by the split rational transport in HD-S10a.

\paragraph{HD-S04 (missing target components).}
For this primitive cyclic construction, all fields other than
\(\mathbb Q(i)\) and \(\mathbb Q(\omega)\), including the other
class-number-one imaginary quadratic fields, are omitted. Explicitly,
the omitted class-number-one fields are
\[
 \mathbb Q(\sqrt{-2}),\quad \mathbb Q(\sqrt{-7}),\quad
 \mathbb Q(\sqrt{-11}),\quad \mathbb Q(\sqrt{-19}),\quad
 \mathbb Q(\sqrt{-43}),\quad \mathbb Q(\sqrt{-67}),\quad
 \mathbb Q(\sqrt{-163}).
\]
The nine-field list is recorded in H.\ M.\ Stark,
\emph{There is no tenth complex quadratic field with class-number one},
Proc.\ Nat.\ Acad.\ Sci.\ U.S.A.\ 57 (1967), 216--221,
p.216, opening paragraph and displayed theorem.
Within the two fields, every ordinary determinant class different from
\([-1]\) is omitted, by the determinant formula just proved.
HD-S03a proves that these are exactly the non-split rational
Hermitian spaces in signature \((3,3)\).
Concrete examples are
\[
 U^2\perp\langle1,-3\rangle\quad(K=\mathbb Q(i)),\qquad
 U^2\perp\langle1,-2\rangle\quad(K=\mathbb Q(\omega)).
\]
They have signature \((3,3)\) and determinants \(-3,-2\).
The primes \(3\) and \(2\), respectively, are inert. A rational field norm
has even valuation at such an inert prime: conjugation fixes the prime,
so the valuation of \(z\bar z\) is twice that of \(z\).
The displayed classes therefore differ from \([-1]\).
Integral Riemann forms are obtained from these Hermitian forms after
clearing a denominator, and an admissible complex structure gives the
corresponding polarized Weil families.
These statements prove the image exclusion of the specified construction;
they make no nonalgebraicity assertion about the omitted targets.

\subsection{Exact period rank}

\paragraph{HD-S05 (the published cubic rank calculation).}
Schoen's \emph{Addendum to: Hodge classes on self-products of a variety
with an automorphism}, Compositio Math. 114 (1998), Lemma 13 and proof,
pp.334--335, proves rank nine for the cubic genus-four source.
Its exact character calculation was rechecked as follows.
For \(G=(\mathbb Z/3)^3\), use inertia vectors
\(e_1,e_2,e_3,-e_1-e_2-e_3\), and the free cyclic subgroup generated
by \(e_1+e_2\).
The two relevant lists of holomorphic characters are
\[
 A=\{(2,2,0),(2,2,1),(2,2,2)\},\quad
 B=\{(0,2,2),(1,1,2),(2,0,2)\}.
\]
Each has multiplicity one under
\(m_k=-1+\frac13\sum_b(k\cdot v_b\bmod3)\).
Their nine pairwise sums are
\[
\begin{pmatrix}
(2,1,2)&(0,0,2)&(1,2,2)\\
(2,1,0)&(0,0,0)&(1,2,0)\\
(2,1,1)&(0,0,1)&(1,2,1)
\end{pmatrix}.
\]
They are distinct; products of two nonzero differentials are nonzero.
Thus the multiplication codifferential has rank exactly nine.
The integer enumeration is recorded in \texttt{certificates.py}.

\paragraph{HD-S06 (the published Gaussian rank calculation).}
Koike, \emph{Algebraicity of some Weil Hodge Classes},
Canad. Math. Bull. 47 (2004), 566--572,
Theorem 2.1, Corollary 2.1 and Remark 2.1, p.568,
prove dominance for the quartic genus-four source.
Its codifferential, equation (1), is
\[
H^0(X,\omega_X\eta)\otimes H^0(X,\omega_X\eta^{-1})
\longrightarrow H^0(X,\omega_X^2).
\]
The source's explicit invariant complete intersections and smooth free-action
open are Proposition 3.1, p.570; the exact determinant calculation is
(6)--(9), pp.571--572.
The supplied computation re-expands all sixteen quadratic identities
and verifies the six-by-six matrix
\[
M(0)=
\begin{pmatrix}
1&0&0&0&0&0\\
0&1&0&0&-2&0\\
0&0&1&0&0&0\\
0&0&0&1&0&0\\
0&0&0&0&-1&0\\
0&0&0&0&0&-1
\end{pmatrix},\qquad\det M(0)=1.
\]
At the second tuple
\((A_1,A_2,A_3;B_1,B_2,B_3;C_1,C_2,C_3)
=(1,0,1;0,1,1;1,1,0)\), it verifies \(\det M=5625\)
and rank six of the seven quadratic forms.
These are polynomial certificates. Their nonempty open conditions meet
the cited smooth free-action open in the irreducible parameter chart;
neither tuple is asserted independently to define a smooth curve.
The remaining three basis vectors in (7) are carried identically,
so the smooth generic rank is nine.

\paragraph{HD-S09 (one exact degeneration certificate for orders 3, 4, 6).}
Let \(X_0\) be the rational nodal genus-four curve with node pairs
\[
(a_j,b_j)=(0,5),(1,7),(2,11),(3,13).
\]
For a root of unity \(\lambda\) of exact order \(d\ge3\), let \(\eta\)
have trivial pullback to the normalization and gluing scalar \(\lambda\)
at each node. Its order is exactly \(d\), since a trivialization on the
connected normalization has only a common constant as automorphism.

The residue rule for \(\omega\eta\) is
\(r_{b_j}=-\lambda r_{a_j}\); the global residue theorem gives
\(\sum_jr_{a_j}=0\). Thus a basis is
\[
f_j\,dz=
\left(\frac1{z-a_j}-\frac1{z-a_4}
-\lambda\left(\frac1{z-b_j}-\frac1{z-b_4}\right)\right)dz,
\quad j=1,2,3.
\]
The inverse-character basis replaces \(\lambda\) by \(\lambda^{-1}\).
Put
\[
D(z)=\prod_{j=1}^4(z-a_j)(z-b_j),\quad
A_j=\frac D{z-a_j}-\frac D{z-a_4},\quad
B_j=\frac D{z-b_j}-\frac D{z-b_4}.
\]
The nine products, with the second basis scaled by \(\lambda\), are
\[
(A_j-\lambda B_j)(\lambda A_k-B_k),\qquad 1\le j,k\le3,
\]
under the common trivialization \((dz)^2/D(z)^2\).
They descend to \(\omega_{X_0}^2\), because the gluing factors cancel.
Riemann--Roch and duality give \(h^0(\omega_{X_0}^2)=9\).

The first nine coefficient rows, corresponding to \(1,z,\ldots,z^8\),
have exact determinant
\[
-C\lambda^4(\lambda-1)^3(\lambda+1)^3
(72\lambda^4+144\lambda^3+419\lambda^2+144\lambda+72),
\]
where
\[
C=2351912680876192340570875868128195756032000000000000000000.
\]
The script \texttt{nodal\_rank\_certificate.py} computes this identity
in \(\mathbb Z[\lambda]\); the JSON contains the complete matrix.
For \(x=\lambda+\lambda^{-1}\in\mathbb R\), the quartic factor divided by
\(\lambda^2\) is \(72(x+1)^2+203>0\).
All other factors are nonzero for \(d\ge3\).
Hence the multiplication map has exact rank nine.

Here is the actual smoothing of the cyclic cover.
Take \(d\) copies of \(\mathbb P^1\), labeled \(r\bmod d\), and identify
\(a_j\) on copy \(r\) with \(b_j\) on copy \(r+1\).
The connected resulting curve has free deck action, \(d\) vertices and
\(4d\) edges in its dual graph, and arithmetic genus \(3d+1\).
At each base node replace \(uv=0\) by \(uv=t_j\), using the same
parameter and local isomorphism at all \(d\) inverse-image nodes.
Equivariant plumbing gives connected etale cyclic covers of smooth
genus-four curves when all \(t_j\ne0\). Compact Riemann surfaces and
their finite holomorphic maps are algebraic.
The eigensheaves stay of exact order \(d\), because the covers stay
connected. Cohomology and base change, using
\(h^0(\eta^{\pm1})=0\), give bundles of ranks \(3,3,9\) for the three
spaces in the multiplication map. Its nonzero determinant persists
at nearby smooth fibers.

The period codifferential is this multiplication, by the character-pair
part of the Kodaira--Spencer contraction and Serre duality.
Exact locators are Koike (2004), (1), p.568, and Lange--Ortega,
\emph{Prym varieties of cyclic coverings}, Geom. Dedicata 150 (2011),
391--403, Proposition 4.1 and Lemma 4.2.
For \(d=3,4,6\), the two primitive characters give the full quadratic
factor, so its unitary target has dimension nine. The computed
nonzero differential minor therefore proves dominance in its actual
polarized lattice component.
Independent modulo-61 reconstruction gave determinants
\(54,8,59\) for primitive roots \(13,11,14\) of orders \(3,4,6\).
This confirms the nonzero rank independently; it is not represented as
a second full symbolic factorization check.

\subsection{A direct effective supplier bound}

\paragraph{HD-S10 (symmetric-product degree bounds).}
For the two named principal source components,
\[
\delta_W(U^3,\mathbb Q(i))\le746496,\qquad
\delta_W(U^3,\mathbb Q(\omega))\le30375.
\]
The Eisenstein bound also holds for the cubic norm Prym with its
restricted Jacobian polarization of type \((1^3,3^3)\).

\paragraph{Proof.}
Let \(m=3\) or \(4\), let \(C\to X\) be as above, and put
\[
N=\{(r_1,\ldots,r_6)\in(\mathbb Z/m)^6:\sum r_i=0\},\quad
G_0=N\rtimes S_6.
\]
The finite maps
\[
C^6\xrightarrow r C^{(6)}
\xrightarrow p W_0=C^6/G_0
\xrightarrow\beta X^{(6)}
\]
have generic degrees \(6!,m^5,m\).
They are free quotients on the locus of six distinct images in \(X\).
Schoen's Theorem 2.0 and the unramified case of Lemma 2.6 supply
\(Q_0\cong\mathbb P^3\) over \(|K_X|\), with nonzero primitive
determinant projection. The map \(\beta\) is etale, by the empty-branch
case of Schoen's Lemma 2.2.
The canonical system has a divisor with distinct zeros, so \(Q_0\)
meets the free locus.

Choose an irreducible component \(Z\subset C^6\) above \(Q_0\), and let
\(W=r(Z)\). If \(H_Z\) is its setwise stabilizer in \(G_0\), generic
freeness gives
\[
\deg(Z/Q_0)=|H_Z|,\quad
\deg(Z/W)=|H_Z\cap S_6|,\quad
e':=\deg(W/Q_0)\le[G_0:S_6]=m^5.
\]
Thus no flatness along the diagonal strata is assumed.
Let \(\xi=[c+C^{(5)}]\) be the ample incidence class on \(C^{(6)}\).
Pullback of a general incidence hyperplane on \(|K_X|\) is \(m\xi\).
Projection formula gives
\[
m^3\int_W\xi^3=e',\qquad \int_W\xi^3\le m^2.                 \tag{S10.1}
\]

We check that this effective component detects.
Let \(U\) be the primitive determinant plane in \(H^6(C^6,\mathbb Q)\).
It is \(G_0\)-invariant, lies in \(H^1(C)^{\otimes6}\), is primitive
for the sum \(H\) of point classes, and has pure type \((3,3)\).
Its cup pairing is nondegenerate by Hodge--Riemann; its group-algebra
projection is self-adjoint.
For \(q_0:C^6\to W_0\), finite-quotient transfer gives
\[
\sum_{g\in G_0}g_*[Z]=\deg(Z/Q_0)\,q_0^*[Q_0].
\]
The right side has nonzero primitive projection, so it pairs
nontrivially with some \(u\in U\).
Invariance yields \(|G_0|\int_Zu\ne0\).
This proves nonzero pairing for the effective component itself.

For \(m=4\), put \(\varphi=(1-\sigma^2)\alpha:C^{(6)}\to B\).
Then \((\varphi r)^*W_K(B)=U\).
Indeed, expand the wedge of six pulled-back primitive one-forms.
Any term with two forms on one curve factor vanishes, because their
character is \(\chi^2\ne1\), while \(H^2(C)\) has trivial character.
The surviving alternating tensor is nonzero, and \(1-\sigma^2\)
acts invertibly on the primitive part.
Projection formula proves that \(\varphi(W)\) has dimension three
and nonzero Weil projection. Its reduced irreducible image is a detector.
For \(m=3\), \(\alpha(W)\) lies in a translate of the norm Prym.
Apply the same argument to \((1-\sigma)\alpha\); because
\(1-\sigma\) is an isogeny on that Prym, the translated Abel image
itself detects.

Write \(\vartheta=\alpha^*\Theta_C\), and let \(\Delta_C\) be the reduced
big diagonal in \(C^{(6)}\). The exact cohomological identity is
\[
\vartheta=(g(C)+5)\xi-\tfrac12[\Delta_C].                    \tag{S10.2}
\]
For clarity, on \(C^6\) the sum-map theta class is \(g(C)H\)
plus the pairwise \(H^1\)-cross terms. The diagonal class on factors
\(i,j\) is \(H_i+H_j\) minus that cross term.
Summing these equalities gives
\[
(\alpha r)^*\Theta_C=(g(C)+5)H-\sum_{i<j}[\Delta_{ij}],
\quad r^*\Delta_C=2\sum_{i<j}\Delta_{ij},
\]
which proves (S10.2), since finite pullback is rationally injective.
A primary source cross-check is
Bastianelli--Kouvidakis--Lopez--Viviani,
\emph{Effective cycles on the symmetric product of a curve, I:
the diagonal cone}, Trans. AMS 372 (2019), 8709--8758,
Example 3.5 (arXiv:1711.07722v3, p.14).
Their cycle-theoretic big diagonal has multiplicity \(4!\) over
the reduced divisor; dividing this explicitly stated map degree
recovers (S10.2).

The component \(W\) is not contained in \(\Delta_C\), because its
generic divisor has six distinct images in \(X\).
Hence \(\Delta_C|_W\) is effective.
Both \(\xi\) and \(\vartheta\) are nef. With \(b=g(C)+5\),
\[
(b\xi)^3-\vartheta^3
=\tfrac12[\Delta_C](b^2\xi^2+b\xi\vartheta+\vartheta^2)
\]
has nonnegative intersection with \(W\); equivalently compute on its
resolution. Thus
\[
\int_W\vartheta^3\le b^3m^2.                               \tag{S10.3}
\]
This uses effectivity of the restricted diagonal, not nefness of that
diagonal.

For \(m=4\), the theorem of the square and
\(\Theta_C|_B=2\Xi\) give
\[
(1-\sigma^2)^*\Xi
=2\Theta_C-\tfrac12(1+\sigma^2)^*\Theta_C\le2\Theta_C.
\]
Mixed nef intersections on \(W\), followed by (S10.3), give
\[
\deg_\Xi\varphi_*W\le8\cdot18^3\cdot16=746496.
\]
For \(m=3\), the translated Abel image has degree at most
\[
15^3\cdot9=30375
\]
for \(\Theta_C|_P\). The explicit principal quotient of HD-S02 has
pullback polarization equal to \(\Theta_C|_P\), so effective pushforward
has the same bound and retains its nonzero Weil projection.
The rank-nine results give a dense source image; the supplied
bounded-Chow closedness extends each bound to its whole named component.

There is also a completely explicit coarse form inequality:
\[
6\sum_i\theta_i-s^*\Theta_C
=\sum_{i<j}(p_i-p_j)^*\Theta_C\ge0
\]
on \(J(C)^6\). Pullback and finite descent yield
\(\vartheta\le6g(C)\xi\).
It independently gives the fallback Gaussian bound \(60742656\).
The sharper value above uses the proved diagonal restriction. \(\square\)

\paragraph{HD-S10a (principal components and fixed polarization transport).}
Over the Eisenstein integers, the matrix with columns
\((1,-\omega)\) and \((1,\bar\omega)\) has determinant \(-1\) and
takes \(U\) to \(I_{1,1}\). Thus the principal cubic quotient is
integrally the balanced product lattice \(I_{3,3}\).
For the Gaussian odd principal lattice, the matrix
\(\begin{psmallmatrix}1&1\\1&-1\end{psmallmatrix}\) in each block
takes \(U\) to \(2I_{1,1}\) and \(I_{1,1}\) to \(2U\).
Three blocks give a degree-64 isogeny with polarization multiplier two.
Hence that explicitly named odd principal component has
\[
\delta_W\le8\cdot746496=5971968.
\]

More generally, fix a split integral maximal-order target of rank six
with polarization elementary-divisor product \(\Delta\).
A primitive isotropic half-lattice extends to an integral basis because
both rings are PIDs. Its Hermitian Gram matrix is
\[
\begin{pmatrix}0&B\\B^*&D\end{pmatrix},\qquad
N(\det B)=|\det H|=\Delta.
\]
The substitutions
\(w'=w-v(B^*)^{-1}D/2,\ u=w'B^{-1}\)
give a hyperbolic basis. The forward and inverse coordinate matrices
are cleared by \(2\Delta\): the only product denominator is
\(2\det B\overline{\det B}\).
The source is already the standard integral \(U^3\), so this single
denominator suffices for the source-to-target comparison.
The resulting isogeny has
\(\phi^*\theta_{\rm target}=(2\Delta)^2\Xi_{\rm source}\).
Consequently
\[
 D_{6,i}(\Delta)=(2\Delta)^6\,746496,\qquad
 D_{6,\omega}(\Delta)=(2\Delta)^6\,30375                      \tag{S10.4}
\]
are uniform effective detector bounds on every such fixed split target.
This is an explicit rational isometry and isogeny construction, not
an identification of distinct integral polarization types.

For a fixed order of conductor \(f\), set
\(\Lambda'=\mathcal O_K\Lambda\). Then \(f\Lambda'\subset\Lambda\)
and \(fE\) is integral on \(\Lambda'\): mixed terms use the integral
\(f\omega\) or \(fi\) action, and the remaining terms use Rosati
compatibility. If \(d=[\Lambda':\Lambda]\), the quotient polarization
satisfies \(q^*\theta'=f\theta\) and
\(\Delta'=f^6\Delta/d\).
Pullback of a three-cycle has degree factor \(d/f^3\).
Since (S10.4) is homogeneous of degree six in \(\Delta\), this yields
\[
\delta_W\le f^{33}D_{6,K}(\Delta).
\]
Use \(f(2+i)\) or \(f(3+\omega)\) in the integral projector.

\paragraph{HD-S07--S08 (retained ambient Bezout audits).}
For comparison, the previously requested ambient calculation is valid
with the full norm Prym.
For the quartic cover its dimension is nine and its type
\((1^6,4^3)\), so
\[
\deg_{3\Theta}W_6=3^6\,13!/7!=900694080,\quad
\deg_{3\Theta}P=3^9\,9!\,4^3=457124290560.
\]
Projection by \(1-\sigma^2\) costs at most \(8/27\) after changing
from \(3\Theta\) to \(\Xi\), giving
\[121993819950101299200.\]
The actual supplier norm fiber is a translate of this nine-dimensional
Prym. Replacing it by the six-dimensional primitive factor in the
intersection has no justified morphism; the degree calculation above
uses the full norm fiber.
For the cubic cover the analogous numbers are
\(110224800\) and \(14171760\), whose product divided by 27 is
\(57854793024000\).
The direct HD-S10 bounds supersede these constants.

\paragraph{HD-S10b (a further published-source corollary).}
The same direct proof in the cubic fourfold case has \(g(C)=7\),
symmetric exponent four, incidence degree at most three, and Abel
coefficient \(7+4-1=10\). It gives
\(\delta_W\le10^2\cdot3=300\) on the principal Eisenstein component,
using Schoen's published rank-four dominance
(Compositio Math.\ 65 (1988), Lemma 3.11, pp.29--30,
and the following proof on p.30).
Its principal quotient is explicitly \(U^2\cong I_{2,2}\) by the
matrix above.

\paragraph{Division, framework, and exact open rows.}
For any of these fixed families, Theorem W of the supplied
HODGE\_DIVISION\_ATTACK R4 converts the displayed integral detector
bound into the explicitly normalized division bound
\(\lceil C_{3,a}(\delta,T)D^4\rceil\).
The separately proved quadratic Gram extraction in this run improves
the dependence on \(D\) to \(D^2\).
No Markman, Severn, or rational-CM-Hodge-conjecture premise is used.

The construction is algebraic: finite quotients, symmetric powers,
Abel--Jacobi maps, and integral isogenies give the actual cycle
morphisms. Their rational homological summands are the primitive
cyclotomic idempotents and the determinant
\(\bigwedge_K^6H^1(B)\), realized in curve powers by the geometric
symmetric projector (the Koszul sign makes it the ordinary exterior
projector on \(H^1\)-tensors). These are the corresponding
lambda-operations in the category of homological motives.
The nodal coefficients are integral; after inverting the covering order,
node differences, and the displayed determinant, the construction spreads
to a finitely generated base over \(\mathbb Z\) by finite presentation
(EGA IV$_3$, Theorem 8.8.2(i)--(ii)).
Its \(R^1\) etale sheaf has weight one at good fibers, and the
determinant has weight six, or zero after twist \((3)\).
Its determinant Euler factor at one K-embedding is
\(1-\det(\operatorname{Frob}|V_\ell)T\);
this records the realization and supplies no extra cycle estimate.

In the project's archive frame the effective uniform detectors establish
condition A, cote 16, scan 050, as corrected in SIGN\_NEXT SN-C03/C04.
The positive Gram-matrix extraction is condition I, scan 049 and
cote 44, Theorem 4. The resultant projector is the correspondence-ring
Cayley--Hamilton inverse of scans 048--049.
These archive locators are carried from the supplied corrected intakes;
no original archive scans were independently inspected in this subtask.
The missing fields and non-split determinant classes remain separate
detector questions: HD-S04 proves the precise source-image exclusion,
and does not replace those questions by nonalgebraicity assertions.

